\documentclass[10pt,a4paper]{amsart}
\usepackage[margin=19mm]{geometry}
\usepackage{amsmath,amssymb,mathtools}
\usepackage[scr=rsfs,cal=euler]{mathalfa}
\usepackage{xfrac}
\usepackage[unicode,hidelinks,bookmarks=false]{hyperref}
\newcommand{\ie}{i.e.\ }
\newcommand{\cf}{cf.\ }
\newcommand{\resp}{respectively\ }
\newcommand{\Subsection}{\subsection}
\providecommand{\nobreakdash}{\nobreak\hbox{-}}
\setlength{\emergencystretch}{2em}
\multlinegap0pt
\usepackage{amssymb}
\usepackage{mathtools}
\usepackage{enumerate}
\usepackage{tikz}
\usepackage{algorithm}\usepackage{algpseudocode}
\newtheorem{theorem}{Theorem}
\title{Towards Real Time Control of Water Engineering with Nonlinear Hyperbolic Partial Differential Equations}
\author{Fabio Difonzo, Michael Holst, Morteza Kimiaei,}
\date{Source: arXiv:2512.14387v4 (CC BY 4.0)}
\begin{document}
\maketitle

\noindent\emph{Source and licence.} Towards Real Time Control of Water Engineering with Nonlinear Hyperbolic Partial Differential Equations. Version arXiv:2512.14387v4 (CC BY 4.0), distributed by arXiv under the Creative Commons Attribution 4.0 International licence, http://creativecommons.org/licenses/by/4.0/, as recorded in the arXiv licence metadata for this exact version and shown on the abstract page https://arxiv.org/abs/2512.14387v4. Full licence notice: \texttt{licenses/arxiv-2512.14387-CC-BY-4.0.txt}.\par\smallskip

\noindent\textit{Editorial scope.} Counted unit: the article's theorem theorem (source0.tex lines 1002--1008). The article's complete original proof of it is delivered with it (source0.tex lines 1009--1011, one \texttt{proof} environment of 3 lines). No mathematics is written, repaired or re-derived: the statement and the proof are the article's own lines, in the article's own notation, and no retained line is substituted or re-flowed.  Retained uncounted as the source-local context the proof needs: lines 948--956; lines 993--999.  Nothing was omitted from any retained span, so the package carries no omission record.  No retained line contains a citation, so the wrapper carries no bibliography and no bibliography entry.  No retained line contains a figure environment, an image-inclusion command or drawing code, so no image is delivered and the package declares no asset dependency.  Editorial formatting only, in addition: an AMS \texttt{amsart} preamble that restates the article's own declarations and macros used by the retained lines, namely the environments \texttt{theorem} on separate counters, together with the standard packages those lines need.  The retained environments print in this document as Theorem 1 in order of appearance, whereas the article prints the same items with the numbering of its own full document.  The complete native source of the article as distributed in the arXiv e-print for this exact version is bundled unchanged as \texttt{sources/191094/source0.tex}.
\begin{theorem}\label{th:shauder}

Let $X$ be a Banach space and either
\begin{itemize}
    \item  $M$ be a nonempty, closed, bounded and convex subset of $X$ and suppose $T:M\to M$ is a compact operator, or
    \item  $M$ be a nonempty, compact,  and convex subset of $X$ and suppose $T:M\to M$ is a continuous operator,
\end{itemize}
then $T$ has a fixed point 
\end{theorem}

\begin{equation}\label{eq:conditionsboundthm}
\begin{array}{l}
K_l\le \frac{l_Q}{u_A}\le \frac{Q^{(0)}}{A^{(0)}} \le \frac{u_Q}{l_A}\le K_u \\
K_l\le \frac{l_Q}{u_A}\le \frac{gA^{(0)}}{2} \le \frac{u_Q}{l_A}\le K_u \\
\tilde{S}_l\le \frac{n_m^2 l^2_Q}{u^{10/3}_A}\le S_f(A^{(0)},Q^{(0)}) = \frac{n_M^2 Q^{(0)}\left\vert Q^{(0)}\right\vert}{(A^{(0)})^{10/3}} \le \frac{n_M^2 u^2_Q}{l^{10/3}_A}\le \tilde{S}_u. \\
\end{array}
\end{equation}

\begin{theorem}
    There exists some $\bar{T}$, such that for all $T\le \bar{T}$, there exists a $\bar{B}>0$ such that there exist $K_l,K_u,\bar{S}_l,\bar{S}_u$ such that with initial conditions satisfying~\eqref{eq:conditionsboundthm}, the following bound holds:
    \begin{equation}
    \max\left\{\|A\|_{L_{\infty}},\left\|A^{-1}\right\|_{L_{\infty}},\|A\|_{L_{\infty}},\left\|A^{-1}\right\|_{L_{\infty}}\right\} \le \bar{B}
    \end{equation}
    for time $0\le t\le T$. As such, the map defined by these iterations is a map between two compact spaces within the set of Colombeau Algebras $\mathcal{C}$, and thus a fixed point exists, which by definition is a solution to the PDE. 
\end{theorem}

\begin{proof}
    The bound follows immediately by integration with respect to $t$ and $x$ and general compactness and boundedness arguments. The result follows as a direct application of Schauder's Fixed Point Theorem~\ref{th:shauder}.
\end{proof}

\end{document}
